% ID: SEL-D0046-OB-E02
% ESTADO: validada. S047.
\textbf{a)} Las ramas son suaves en sus interiores. En $x=2$ igualamos el valor de la rama izquierda y el límite de la derecha: $a-6b=30$. Para la derivabilidad, una vez establecida la continuidad, igualamos las derivadas laterales: $-5b=-60/4=-15$. Así $\boxed{b=3,\ a=48}$.

\textbf{b)} En la rama izquierda, $f'(x)=-6x-3$ cambia de positivo a negativo en $x=-1/2$. En la rama derecha, $f'(x)=-60/x^2<0$, y en la unión $f'(2)=-15$. Por tanto, $f$ crece en $(-\infty,-1/2)$ y decrece en $(-1/2,\infty)$. Tiene un máximo absoluto en $x=-1/2$, de valor $\boxed{195/4}$. No tiene mínimo, pues $f(x)\to-\infty$ cuando $x\to-\infty$.
